%
\documentclass[a4paper, 10pt, conference]{ieeeconf}

\IEEEoverridecommandlockouts{}
\overrideIEEEmargins{}

\usepackage{graphics, epsfig, url}
\usepackage{amsmath, amssymb, mathptmx}
\usepackage[colorlinks=true]{hyperref}

\title{
    \LARGE \bf
    Task-parameterised Gaussian Processes
    \\ [0.5em]
    \large MTech Project MidTerm Report
}
\author{
    Ashrith Sagar Yedlapalli$^{1}$
    \thanks{$^{1}$ ashrithy@iisc.ac.in
    {\tt\small ashrithy@iisc.ac.in}}
}

\begin{document}

\maketitle
\thispagestyle{empty}
\pagestyle{empty}

\begin{abstract}
    Task-parameterised Gaussian Processes
\end{abstract}

\section{INTRODUCTION}

Robot learning from demonstrations (LfD) has emerged as a powerful paradigm for teaching robots complex tasks without explicit programming.
In interactive LfD, non-expert users can kinesthetically guide the robot to perform tasks like pick-and-place, dressing, or cleaning.
However, a major challenge is generalizing these learned policies to novel situations, such as varying object positions, arm postures, or surface geometries.

\section{RELATED WORK}

Movement primitives (MPs) provide compact, reusable representations of robotic motions, enabling efficient imitation learning, generalization, and adaptation.
Dynamic Movement Primitives (DMPs), introduced by Ijspeert et al.~\cite{Ijspeert2013}, model both discrete and rhythmic behaviors as autonomous dynamical systems.
A DMP consists of a canonical linear system generating a phase variable and a transformation system (second-order spring-damper) driven by a learned nonlinear forcing term approximated via locally weighted regression on basis functions.
Demonstrations are encoded into weights, allowing straightforward modulation of amplitude, goal, duration, and coupling with external forces or other primitives—making DMPs a cornerstone for trajectory generation in imitation learning.

Probabilistic Movement Primitives (ProMPs) by Paraschos et al.~\cite{10.5555/2999792.2999904} extend DMPs into a fully probabilistic framework.
Trajectories are represented as a linear combination of radial basis functions whose weights follow a multivariate Gaussian distribution learned from multiple demonstrations.
This captures movement variability, supports stochastic execution, and enables conditioning on via-points or final states through closed-form Bayesian updates.
Kernelized Movement Primitives (KMPs), proposed by Huang et al.~\cite{Huang2019}, replace fixed basis functions with kernel ridge regression, yielding a non-parametric representation.
KMPs naturally handle multiple via-point constraints, obstacle avoidance, and fusion of primitives without explicit phase alignment, offering superior flexibility for complex task adaptation in interactive imitation scenarios.

\section{METHODOLOGY}

\section{RESULTS \& DISCUSSIONS}

\section{CONCLUSIONS}

\addtolength{\textheight}{-12cm}

\section*{APPENDIX}

\section*{ACKNOWLEDGMENT}

\bibliographystyle{IEEEtran}
\bibliography{references}

\end{document}
